\section{Related Work}
\label{sec:related-work}

\noindent
\textbf{Self-Supervised Learning (SSL)} leverages unlabeled data to learn meaningful representations for downstream tasks. In computer vision, three main approaches are: (1) masked prediction \cite{he2022masked}, (2) contrastive learning like MoCo \cite{he2020momentum} and SimCLR \cite{chen2020simple}, and (3) self-distillation methods such as DINO \cite{DINO}. These techniques have enabled large-scale training of models with emergent abilities \cite{wei2022emergent}.

\noindent
\textbf{Training on Imbalanced Datasets} often leads to inferior performance and bias toward majority classes \cite{johnson2019survey}. While SSL methods are more robust to imbalance than supervised approaches \cite{liu2021self}, they still exhibit diminished performance on imbalanced data. Current solutions are primarily model-centric, attempting to learn arbitrary feature priors \cite{assran2210hidden}. Data-centric approaches that directly complement imbalanced datasets remain underexplored.

\noindent
\textbf{Out-Of-Distribution (OOD) Detection} is crucial for enhancing model robustness by maintaining high-quality datasets \cite{whang2023data}. Distance-based approaches \cite{sun2022out} identify OOD samples by measuring their distance from in-distribution samples in the embedding space. These non-parametric methods offer flexibility without distributional assumptions, making them suitable for imbalanced, unlabeled data.

\noindent
\textbf{Diffusion-Based Image Generation} models systematically add and then remove noise to generate high-quality images \cite{ho2020denoising}. Stable Diffusion \cite{rombach2022highresolution}, particularly Stable Diffusion 2 UnClip, leverages CLIP embeddings to generate semantically similar images to the input context. Its ability to perform image-to-image generation while preserving semantic content makes it well-suited for augmenting self-supervised learning datasets.